%==============================================================================
%  MOTIVATION LETTER -- university portal or graduate programme
%  Register: narrative is WANTED here, unlike the cover letter. 500-650 words,
%  one page if it fits, two at most. Respect any word limit the portal states.
%
%  The key difference from cover-example.tex: a motivation letter DOES want
%  "why this group" -- but specific (their method line, their cohort, a named
%  paper), never "your beautiful campus and vibrant city".
%==============================================================================
\documentclass[11pt,a4paper]{article}
\usepackage{ssletter}
\begin{document}

\letterhead
\lettermeta{\fillme{date}}
\recipient{\fillme{Admissions / Selection Committee, Department, University}}

\opening{Dear \fillme{Members of the Selection Committee / Dr Surname}}

I am applying to \fillme{programme or position, with reference number} at
\fillme{institution}. I am an electrical engineer working on causal inference for
resting-state fMRI, and I want to spend a PhD on a problem I have already run into from the
applied side: \fillme{the problem, in one clause}.

% ---- Where you come from, methodologically. Narrative is appropriate here.
My route into this was through control systems rather than neuroscience, which shaped how I
approach it. During my MSc at the University of Tehran I applied constraint-based (PC),
score-based (GES) and functional (LiNGAM) causal discovery to resting-state fMRI in
schizophrenia (\fillme{n} patients, \fillme{n} controls, \fillme{atlas}), combining the
recovered directed graphs with graph node embeddings to classify patients
(\fillme{result vs.\ the correlation-based baseline}). I then joined a clinical group at
Tehran University of Medical Sciences to estimate effective connectivity in an epilepsy cohort
(\fillme{n}), where working alongside radiologists taught me the distance between a
statistically significant connectivity result and a claim a clinician will act on. Both
studies met the same wall — small, imbalanced, single-site cohorts — which is why my current
manuscript builds generative augmentation for connectivity data rather than another
classifier.

% ---- Why THIS group. Be specific: name the method line or a paper.
\fillme{Why this group specifically: name their method line, cohort, or facility, and one
paper you have actually read. One or two sentences on what they do that you cannot do
elsewhere. Never mention the city.}

% ---- What you would work on. A concrete direction, not "I am eager to learn".
Concretely, I would want to work on \fillme{your proposed direction, in one or two
sentences}. \fillme{Why it is open: what is unresolved, and why their setting is the right
place to attempt it.}

% ---- What you bring that the requirements did not ask for.
Beyond the stated requirements I bring two things. First, engineering scale: I currently build
production machine-learning pipelines over 70 million records and ran my fMRI analyses on HPC
infrastructure, so large simulation studies and multi-site pipelines are work I have done
rather than would be learning. Second, \fillme{second differentiator: a complementary methods
family, clinical collaboration, or graduate teaching}.

I was ranked 4th nationally among approximately 20,000 candidates in the Iranian MSc entrance
examination and held full scholarships for both degrees. My publications, code and figures are
at \fillme{website}.

I can begin \fillme{date} and am ready to relocate. Thank you for considering my application;
I would welcome the opportunity to discuss the project.

\closing
\end{document}
